\section{Objetivos de la propuesta}

\noindent\textbf{Objetivo general.} Diseñar, construir y validar experimentalmente una máquina automática de doblado por línea térmica para láminas de acrílico de hasta \SI{25}{\centi\meter} de longitud de doblez y espesores de \SIrange{3}{6}{\milli\meter}, con control realimentado de temperatura e interfaz táctil para seleccionar ángulo y posición, destinada al Laboratorio de Mecatrónica de la Universidad EAFIT.

\noindent\textbf{Objetivos específicos.}
\begin{enumerate}
\item Caracterizar la ventana de conformado del PMMA para espesores de \num{3}, \num{4.5} y \SI{6}{\milli\meter}, obteniendo curvas de temperatura contra tiempo, ancho de franja plastificada y \emph{springback} frente a temperatura y tiempo de sostenimiento.
\item Diseñar el subsistema térmico (resistencia tubular, canal reflector, aislamiento y sensado) con un modelo de transferencia de calor validado con termopares, de modo que la franja se mantenga dentro de \SI{\pm 5}{\celsius} en los \SI{25}{\centi\meter}.
\item Diseñar e implementar el subsistema electromecánico de sujeción, posicionamiento longitudinal y plegado motorizado, con resolución angular y de posición iguales o mejores que \SI{0.5}{\degree} y \SI{0.5}{\milli\meter}.
\item Implementar en Arduino el control: lazo PID de temperatura con relé de estado sólido, secuencia de plegado con compensación de \emph{springback} e interfaz táctil de selección de espesor, ángulo y posición.
\item Integrar, construir y poner en marcha el prototipo, con los elementos de seguridad eléctrica y térmica exigidos en un laboratorio docente.
\item Validar el desempeño con al menos \num{30} dobleces por condición, verificando error angular no mayor a \SI{\pm 2}{\degree}, error de posición no mayor a \SI{\pm 1}{\milli\meter}, repetibilidad y ausencia de defectos ópticos.
\end{enumerate}
